% !TeX program = pdflatex
\documentclass[11pt,a4paper]{article}
\usepackage[T1]{fontenc}
\usepackage[utf8]{inputenc}
\usepackage{lmodern,amsmath,amssymb,amsthm,mathtools,mathrsfs,microtype}
\usepackage[margin=24mm]{geometry}
\usepackage[hidelinks]{hyperref}
\usepackage{xurl}
\setlength{\emergencystretch}{2em}
\allowdisplaybreaks[2]
\newtheorem{theorem}{Theorem}[section]
\newtheorem{lemma}[theorem]{Lemma}
\newtheorem{corollary}[theorem]{Corollary}
\newtheorem{proposition}[theorem]{Proposition}
\theoremstyle{remark}\newtheorem{remark}[theorem]{Remark}
\newcommand{\Ric}{\operatorname{Ric}}
\newcommand{\tr}{\operatorname{tr}}
\newcommand{\diver}{\operatorname{div}}
\newcommand{\im}{\operatorname{im}}
\newcommand{\Hharm}{\mathcal H}
\newcommand{\GP}{\mathcal G_P}
\newcommand{\ip}[2]{\langle #1,#2\rangle}
\numberwithin{equation}{section}
\title{Harmonic rank alignment and negative-scalar\\[3pt]trace-free Ricci criticality}
\author{Qiushi research continuation}
\date{8 October 2026}
\begin{document}
\maketitle
\begin{abstract}
For the unrestricted functional $P(g)=\int|\Ric_g^\circ|^2$ on a closed
oriented four-manifold, we exclude a constant-negative-scalar non-Einstein
critical metric whose mixed curvature block has image in a harmonic
anti-self-dual line. Neither constant amplitude nor removability of the other
factor is assumed. The proof retains the full pointwise Euler--Lagrange
operator, shows that the Ricci tensor must vanish at zeros of the aligned
harmonic form, and tests the remaining equation at a maximum of its squared
norm. A consequence is that a negative-scalar Riemannian Einstein--Maxwell
metric can be $P$-critical only if it is Einstein. This strengthens the
harmonic-alignment exclusions in the supplied manuscripts; it does not prove
smooth attainment of the unrestricted infimum on the Godeaux blowup.
\end{abstract}

\section{Functional, curvature block, and the statement}
Let $(M^4,g)$ be smooth, closed, connected and oriented. Set
\[
 E=\Ric-\frac{s}{4}g,\qquad P(g)=\int_M|E|^2\,dV_g.
\]
Two-forms have the form norm, with $|e^{12}+e^{34}|^2=2$.
We compose their skew endomorphisms using $g$ and identify the resulting
self-adjoint endomorphisms with symmetric two-tensors. Write
\[
 \mathcal R=\begin{pmatrix}W^++sI/12&B^*\\B&W^-+sI/12\end{pmatrix},
 \qquad B:\Lambda^+\longrightarrow\Lambda^-,\qquad |E|^2=4|B|^2.
\]
The curvature contraction is $\mathscr R(H)_{ij}=R_{ikjl}H^{kl}$,
with $R_{ikjl}g^{kl}=\Ric_{ij}$ and positive sectional curvature on a sphere.
The rough Laplacian $L=\nabla^*\nabla$ is nonnegative.

\begin{theorem}\label{thm:main}
Suppose $s$ is a negative constant and $0\ne h\in\Hharm^-(M,g)$ satisfies
\begin{equation}\label{eq:align}
 \im B_x\subset\mathbb R h_x\qquad(x\notin Z(h)).
\end{equation}
Let $\GP$ be the unrestricted $L^2$ gradient of $P$. Then
\begin{equation}\label{eq:infty}
 \boxed{\|\GP\|_\infty\ge\frac{|s|}{2}\|E\|_\infty.}
\end{equation}
If $E\not\equiv0$, there is a smooth nonnegative cutoff $\chi$, supported
away from $Z(h)$, such that the volume-form-preserving metrics
\begin{equation}\label{eq:path}
 g_t(X,Y)=g\bigl(e^{-t\chi E^\sharp}X,Y\bigr)
\end{equation}
satisfy $P(g_t)<P(g)$ for all sufficiently small $t>0$.
In particular, a $P$-critical metric satisfying \eqref{eq:align} is Einstein.
\end{theorem}

The condition is exact rank alignment with a specified global harmonic
form. It is not the assertion that a general $P$-minimizer has this structure.
No lower bound for $|h|$ is required. The proof will not extend or integrate
an unbounded quotient across $Z(h)$.

\section{The pointwise first variation}
The classical gradients of $\int|\Ric|^2$ and $\int s^2$ in
\cite[(2.5), (2.9)]{GV}, combined as $P=\int|\Ric|^2-\frac14\int s^2$,
give at constant scalar curvature
\begin{equation}\label{eq:grad}
 \GP=LE-2\mathscr R(E)+\frac12|E|^2g,
 \qquad DP_g[k]=\int_M\ip{\GP}{k}\,dV_g.
\end{equation}
The trace of the full gradient at an arbitrary metric is a nonzero multiple
of $\Delta s$; hence every smooth closed $P$-critical metric has constant
scalar curvature. Scale invariance makes unit-volume criticality equivalent
to unrestricted criticality.

\begin{lemma}[Opposite-duality tensor algebra]\label{lem:algebra}
If $E=\phi\circ\psi$ with $\phi\in\Lambda^+$ and $\psi\in\Lambda^-$,
then
\begin{align}
 |E|^2&=|\phi|^2|\psi|^2,&
 (E^\sharp)^2&=\frac{|E|^2}{4}I,\label{eq:norm}\\
 \mathscr R(E)&=(W^+\phi)\circ\psi+\phi\circ(W^-\psi)
               -\frac{s}{12}E+\frac{|E|^2}{4}g.
 \label{eq:action}
\end{align}
Here $E$ is the trace-free Ricci tensor of the curvature used in
\eqref{eq:action}.
\end{lemma}
\begin{proof}
Opposite-duality skew endomorphisms commute and their product realizes the
isometric identification $\Lambda^+\otimes\Lambda^-\simeq S_0^2$:
\begin{equation}\label{eq:inner}
 \ip{\alpha\circ\beta}{\gamma\circ\delta}
 =\ip{\alpha}{\gamma}\ip{\beta}{\delta}
 \quad(\alpha,\gamma\in\Lambda^+,\ \beta,\delta\in\Lambda^-).
\end{equation}
At a point with both forms nonzero one can choose an oriented orthonormal
frame with $\phi=a(e^{12}+e^{34})$, $\psi=b(e^{12}-e^{34})$.
Their product is $ab\operatorname{diag}(-1,-1,1,1)$, proving
\eqref{eq:norm}; zero cases follow by continuity.

For completeness decompose the curvature tensor as
\[
 R_{ikjl}=W_{ikjl}+\frac12
 (E_{ij}g_{kl}+E_{kl}g_{ij}-E_{il}g_{kj}-E_{kj}g_{il})
 +\frac{s}{12}(g_{ij}g_{kl}-g_{il}g_{kj}).
\]
Contraction with $E^{kl}$ gives
\[
 \mathscr R(E)=W(E)+\frac12|E|^2g-E^2-\frac{s}{12}E.
\]
Under the same product identification,
$W(E)=(W^+\phi)\circ\psi+\phi\circ(W^-\psi)$.
This is verified by multiplication in the displayed frame for arbitrary
symmetric trace-free $3\times3$ matrices $W^+,W^-$; equivariance gives the
identity in all frames. Equation \eqref{eq:norm} proves \eqref{eq:action}.
\end{proof}

\begin{proposition}[Local harmonic-product formula]\label{prop:local}
On any open set where $E=\phi\circ\psi$ and both opposite-duality forms
are smooth and harmonic, constant scalar curvature implies
\begin{equation}\label{eq:pointgrad}
 \boxed{\GP=-\frac{s}{2}E-
 2\sum_{i=1}^4(\nabla_{e_i}\phi)\circ(\nabla_{e_i}\psi).}
\end{equation}
Writing $a=|\phi|^2$ and $b=|\psi|^2$, one has pointwise
\begin{equation}\label{eq:contract}
 \boxed{\ip{\GP}{E}=-\frac{s}{2}ab-\frac12\ip{da}{db}.}
\end{equation}
\end{proposition}
\begin{proof}
The harmonic two-form Weitzenb\"ock equations are
\[
 L\phi=2W^+\phi-\frac{s}{3}\phi,\qquad
 L\psi=2W^-\psi-\frac{s}{3}\psi.
\]
Apply the product rule
\[
 L(\phi\circ\psi)=(L\phi)\circ\psi+\phi\circ(L\psi)
 -2\sum_i(\nabla_{e_i}\phi)\circ(\nabla_{e_i}\psi).
\]
Substitution of \eqref{eq:action} into \eqref{eq:grad} cancels both Weyl
terms and the scalar multiple of $g$, leaving \eqref{eq:pointgrad}.
Equation \eqref{eq:inner} then gives
\[
 \sum_i\ip{(\nabla_{e_i}\phi)\circ(\nabla_{e_i}\psi)}{\phi\circ\psi}
 =\frac14\ip{da}{db},
\]
which proves \eqref{eq:contract}. This is a pointwise calculation, so there
are no discarded boundary or zero-set terms.
\end{proof}

\section{Harmonic zeros and the maximum argument}

\begin{lemma}[Alignment at a zero]\label{lem:zeros}
Let $h\not\equiv0$ be a harmonic anti-self-dual two-form on a connected
four-manifold, and let $B:\Lambda^+\to\Lambda^-$ be smooth and satisfy
\eqref{eq:align}. Then $B=0$ at every point of $Z(h)$.
\end{lemma}
\begin{proof}
The zero set of a nonzero harmonic form has empty interior by unique
continuation for the Hodge--Dirac operator \cite{Baer}. Thus $\operatorname{rank}B\le1$
everywhere, since its $2\times2$ minors vanish on a dense set.
Suppose $h(p)=0$ and $B(p)\ne0$. After shrinking to a connected ball $U$
choose a smooth local section $v$ of $\Lambda^+$ with $Bv\ne0$ throughout
$U$, and put $\theta=Bv/|Bv|$. This is a smooth unit anti-self-dual form.
The alignment gives $h=f\theta$ on $U$, including its zero set, for the
smooth real function $f=\ip{h}{\theta}$.
Since $dh=0$,
\[
 df\wedge\theta+f\,d\theta=0.
\]
A nonzero anti-self-dual two-form is nondegenerate. Therefore
$\alpha\mapsto\alpha\wedge\theta$ is an isomorphism from one-forms to
three-forms; there is a smooth one-form $\omega$ on $U$ with $df=f\omega$.
Along every path from $p$, this is a scalar homogeneous linear ODE.
The initial value $f(p)=0$ implies $f=0$ on $U$, contradicting unique
continuation of $h$. Hence $B(p)=0$.
\end{proof}

\begin{proof}[Proof of Theorem \ref{thm:main}]
On $M\setminus Z(h)$ set
\[
 \beta_h=\frac{B^*h}{|h|^2},\qquad \phi=-2\beta_h.
\]
The mixed-block/tensor identification gives $E=\phi\circ h$.
Contracted Bianchi at constant scalar curvature gives $\diver E=0$.
For a self-dual $\phi$ and closed anti-self-dual $h$,
\begin{equation}\label{eq:div}
 \diver(\phi\circ h)=h(d^*\phi).
\end{equation}
To check the only potentially extra term, contract $dh=0$ with
$\phi^{ik}$. The two differentiated terms give the same contraction,
while $\phi^{ik}\nabla_jh_{ik}=0$ by opposite duality. Thus the derivative
falling on $h$ vanishes. This proves \eqref{eq:div} with the indicated
skew-endomorphism convention. As $h$ is invertible off its zeros,
$d^*\phi=0$, and self-duality gives $d\phi=0$. The local product formula
therefore applies without any assertion of a global smooth extension of
$\phi$.

Assume $E\not\equiv0$ and choose $p$ where $|E|^2$ has its positive
maximum. Lemma \ref{lem:zeros} and $|E|^2=4|B|^2$ imply $h(p)\ne0$.
Hence both $a=|\phi|^2$ and $b=|h|^2$ are positive near $p$ and
$|E|^2=ab$. At $p$,
\[
 0=d(ab)=b\,da+a\,db,
 \qquad \ip{da}{db}=-\frac ba|da|^2\le0.
\]
Equation \eqref{eq:contract} now gives
\begin{equation}\label{eq:max}
 \ip{\GP}{E}(p)=\frac{|s|}{2}|E|^2(p)
                  +\frac{b(p)}{2a(p)}|da|^2(p)>0.
\end{equation}
Cauchy--Schwarz at $p$ proves \eqref{eq:infty}. By continuity a nonnegative
smooth cutoff $\chi\not\equiv0$ can be chosen near $p$ so that
\[
 DP_g[-\chi E]=-\int_M\chi\ip{\GP}{E}\,dV_g<0.
\]
The exponential in \eqref{eq:path} is positive and has determinant one
because $\tr E=0$. It therefore gives the claimed actual metrics and
first-order strict descent. For $E=0$ all assertions are immediate.
\end{proof}

\begin{remark}[The constant is sharp]
For a product of two closed constant-negative-curvature surfaces with
unequal Gaussian curvatures, the Ricci tensor is parallel and the two
factors in its harmonic-product expression can both be chosen parallel.
Then \eqref{eq:pointgrad} becomes $\GP=-sE/2$, and equality holds in
\eqref{eq:infty}. The theorem excludes criticality, not existence of such
Einstein--Maxwell metrics.
\end{remark}

\section{Einstein--Maxwell and the Q3 interface}

\begin{corollary}\label{cor:EM}
Let $(g,F)$ be a smooth Riemannian Einstein--Maxwell solution on a closed
oriented four-manifold, in the convention
\[
 dF=d^*F=0,\qquad (\Ric+F\circ F)_0=0.
\]
If its scalar curvature is negative, then $P$-criticality of $g$ is
equivalent to $g$ being Einstein. Every non-Einstein member admits the
localized volume-preserving descent in \eqref{eq:path}.
\end{corollary}
\begin{proof}
The Maxwell stress tensor is divergence-free, so contracted Bianchi makes
$s$ constant. The standard four-dimensional decomposition \cite{LeBrunEM}
is $(F\circ F)_0=2F^+\circ F^-$. Hence
\[
 E=(-2F^+)\circ F^-.
\]
Both factors are globally smooth harmonic forms. If $F^-\equiv0$, the
metric is already Einstein. Otherwise its mixed block is aligned with
$F^-$, and Theorem \ref{thm:main} applies. Conversely, an Einstein metric
has $P=0$ and is an absolute $P$-minimum.
\end{proof}

\begin{corollary}[Strict rank remainder at a putative minimizer]
Suppose $g_*$ is a smooth non-Einstein $P$-critical metric with $s<0$.
For every $0\ne h\in\Hharm^-(M,g_*)$, define off $Z(h)$
\[
 B_{\perp,h}=B-\frac{h\otimes B^*h}{|h|^2}.
\]
Then $\int|B_{\perp,h}|^2>0$. If $b^-(M)>0$, one further has
\begin{equation}\label{eq:rankgap}
 \min_{\substack{h\in\Hharm^-(g_*)\\\|h\|_2=1}}
       \int_M|B_{\perp,h}|^2\,dV_{g_*}>0.
\end{equation}
The positive number in \eqref{eq:rankgap} may depend on $g_*$.
\end{corollary}
\begin{proof}
A vanishing integral would give \eqref{eq:align} on the complement of
$Z(h)$ by smoothness there, contradicting Theorem \ref{thm:main}.
The unit sphere of the harmonic space is compact, and convergence in it
is smooth. A nonzero harmonic form has a measure-zero nodal set
\cite{Baer}; dominated convergence, using $|B_{\perp,h}|^2\le|B|^2$,
shows that the displayed integral is continuous in $h$. Its positive
minimum proves \eqref{eq:rankgap}.
\end{proof}

\paragraph{What this changes in the supplied proof.}
The supplied continuation \cite{Project} separates full flux equality from
weaker rank alignment. Its global Einstein--Maxwell conclusion required
boundedness of the quotient factor at zeros of the harmonic form. Theorem
\ref{thm:main} removes that intermediate requirement for the negative-scalar
criticality question: it uses the original smooth mixed curvature block to
show $E=0$ at those zeros and works at a positive maximum away from them.
Both amplitudes are allowed to vary. The earlier integrated formula
\[
 DP_g[E]=-\frac{s}{2}P-\frac12\int
       \ip{d|\phi|^2}{d|\psi|^2}
\]
is recovered when both factors are global, but its unsigned correlation is
not used to obtain the contradiction. The full pointwise equation is the
additional information.

For the fixed Godeaux blowup the supplied lower and strict upper bounds
remain $16\pi^2/3\le P_*(M)<8\pi^2$. No argument here proves smooth
attainment of $P_*(M)$ or forces an arbitrary minimizing sequence into
harmonic rank alignment. The zero-mode ALE and partial-collapse problems
are not removed by this candidate-class exclusion. The present conclusions
use classical curvature identities and unique continuation, not FD2 or the
project's conditional gluing expansion.

\begin{thebibliography}{9}
\bibitem{GV} M. J. Gursky and J. A. Viaclovsky,
\emph{Rigidity and stability of Einstein metrics for quadratic curvature
functionals}, arXiv:1105.4648v2 (2013).
\url{https://arxiv.org/abs/1105.4648}. Equations (2.5), (2.9).
\bibitem{LeBrunEM} C. LeBrun,
\emph{The Einstein--Maxwell equations, extremal K\"ahler metrics, and
Seiberg--Witten theory}, arXiv:0803.3734v2 (2008).
\url{https://arxiv.org/abs/0803.3734}.
\bibitem{Baer} C. B\"ar,
\emph{Zero sets of solutions to semilinear elliptic systems of first order},
Invent. Math. 138 (1999), 183--202.
\url{https://arxiv.org/abs/math/9805065}.
\bibitem{Project} Supplied project manuscript,
\emph{Einstein Gaps on Four-Manifolds: Subsequent Mathematical Developments},
8 October 2026, subsection ``The weaker rank-aligned system'';
and the preceding note \emph{A genuine descent direction on the
harmonic-alignment branch}. These are the research context, not additional
analytic hypotheses for Theorem \ref{thm:main}.
\end{thebibliography}
\end{document}
